\section{Introduction}
\label{sec:intro}

A vision-language model rarely sees an image alone. It is shown a photo with a caption scraped from the web, a page returned by a retrieval system, or a user's statement about what the photo contains, and any of these can be wrong. What a model should do is clear: notice that the text and the image disagree, answer from the image when the question is about the image, and be able to say which input is wrong and how. Each of these abilities has been studied, but mostly on its own. Hallucination benchmarks ask whether a model describes objects that are not there~\cite{rohrbach2018chair,li2023pope,guan2024hallusionbench}; foil and compositionality benchmarks ask whether it can tell a correct caption from a minimally altered one~\cite{shekhar2017foil,parcalabescu2022valse,thrush2022winoground,hsieh2023sugarcrepe}; and recent work shows that VLMs often follow misleading text over the image~\cite{deng2025words}. Whether a model that notices a conflict also acts on it, and whether it can describe what it noticed, has not been measured on the same items.

This matters because the three abilities support different uses. A system that flags inconsistent inputs needs detection; a system that answers questions needs resolution; a system that explains or repairs its inputs needs attribution. If they come apart, good performance on one says little about the others, and a model that reports a conflict can still act as if it had not seen one.

We introduce \name, a benchmark that measures detection, resolution and attribution jointly on controlled conflicts (\cref{fig:method}). Each record pairs a real photo with a short description built from its GQA scene graph~\cite{hudson2019gqa}, in which exactly one span is changed: an object, a colour or a spatial relation. Edits are graded into three levels of perceptual difficulty, and faithful descriptions of the same kind serve as controls for false alarms. Every model answers the same records under ten probes: a yes/no detection question, a multiple-choice question about the edited fact, a structured error report, and controls that measure perception from the image alone, reading from the text alone, the effect of an instruction to rely on the image, and the effect of conversational turn structure.

We evaluate eleven open models from four families, including a scaling ladder within one family from 2B to 32B parameters, and find that the abilities come apart in consistent ways:
\begin{itemize}[leftmargin=*,itemsep=1pt,topsep=2pt]
\item \textbf{Noticing is not acting.} Models flag 72--98\% of conflicts, yet for six of eleven the most common outcome is a flagged conflict answered without the image. Among conflicts a model perceives correctly from the image alone \emph{and} flags, it still answers from the text 26--88\% of the time (\cref{fig:teaser}).
\item \textbf{Scale does not fix the default.} Answering from the image improves up to 7--8B in every family, but both 32B models fall back by 11 and 16 points relative to their 7--8B siblings, although told to rely on the image they do so on 86\% and 80\% of conflicts. False alarms on faithful descriptions rise from the 2--3B to the 7--8B model in every family.
\item \textbf{Self-reports are unreliable.} Asking a model to report the error before answering helps, but control probes show that for most models the turn structure contributes more than the report. Models locate the edited words well yet rarely say which input is wrong, and five models state ``no error'' in most of their reports on conflicts they flag in the detection probe.
\end{itemize}

Our contributions are: (1) \name, a benchmark of single-span, severity-graded image--text conflicts with faithful controls\pending{, validated by two annotators}; (2) a probing protocol that measures detection, resolution and attribution on the same records and separates perception, arbitration, instruction following and turn structure; and (3) an evaluation of eleven open VLMs with image-clustered bootstrap intervals and paired tests, showing that the three abilities come apart and that scale changes the default behaviour in a non-monotonic way. All code, data and model outputs will be released.
